\documentclass[11pt]{article}

\usepackage[a4paper,margin=1in]{geometry}
\usepackage{amsmath,amssymb,amsthm,mathtools}
\usepackage{booktabs}
\usepackage{graphicx}
\usepackage{xcolor}
\usepackage[colorlinks=true,linkcolor=blue!55!black,citecolor=blue!55!black,urlcolor=blue!55!black]{hyperref}
\usepackage{microtype}

\newtheorem{theorem}{Theorem}
\newtheorem{proposition}{Proposition}
\newtheorem{remark}{Remark}
\newcommand{\R}{\mathbb{R}}
\newcommand{\BV}{BV_{\mathrm{loc}}}
\newcommand{\HH}{\mathcal{H}}
\newcommand{\V}{\mathsf{V}}
\newcommand{\Tmap}{\mathcal{T}}
\DeclareMathOperator{\diag}{diag}

\title{A Jump-Line Obstruction to Fixed-Point\\Preservation of Variability}
\author{Automatic Math Discovery Run \texttt{fa\_banach\_001}\\Model: GPT5.6}
\date{26 August 2026}

\begin{document}
\maketitle

\begin{center}
\fcolorbox{blue!45!black}{blue!4}{\parbox{0.9\textwidth}{
\textbf{Status.} Candidate substantial partial result, likely valid.  This is
an obstruction to a natural fixed-point invariance principle, not a full
resolution of the source paper's open-ended request for some suitable
assumptions.
}}
\end{center}

\begin{abstract}
Hinz, T\"olle, and Viitasaari ask for assumptions under which the integral
process associated with a $BV$ coefficient is itself variable, so that a
standard fixed-point construction might replace their Doss transform.  We
show that the most immediate invariance statement fails in dimension two under
strong baseline hypotheses.  A bounded, symmetric, uniformly positive
definite, piecewise-constant coefficient can satisfy their determinant,
curl-free-inverse, and angle conditions, while a smooth driver with both
coordinates active sends an $(s,p)$-variable input path exactly onto the
coefficient's jump line.  The output is then not variable for any positive
variability exponent.  The example works for every prescribed finite $p$.
It shows that a successful fixed-point theorem needs an output-side
transversality or occupation-energy hypothesis, not only coefficient
nondegeneracy and input variability.  A short scalar positive result is added
to illustrate the missing mechanism.
\end{abstract}

\section{Source problem and scope}

The source is M. Hinz, J. M. T\"olle, and L. Viitasaari,
\emph{Variability of paths and differential equations with
$BV$-coefficients}~\cite{HTV}.  In Section 1.2, PDF page 6, after describing
their Doss-transform existence theory, the authors state that the main open
problem is to prove, under reasonable assumptions, that ``the integral
process itself will be variable.''  Their goal is to make standard fixed-point
arguments available.

\begin{figure}[ht]
  \centering
  \includegraphics[width=\textwidth]{figures/open_problem_crop.png}
  \caption{The source statement on PDF page 6.  The image is a direct crop of
  the source paper, not a retyped substitute.}
\end{figure}

Because ``reasonable assumptions'' are deliberately unspecified, no example
can refute the question as written.  The result below instead tests a strong
and natural baseline: bounded $BV$ coefficients satisfying all three
structural assumptions used by the source's multidimensional Doss theory,
smooth drivers, and variable input paths.  It proves that these conditions do
not make the raw Picard map invariant on variability classes.

\section{Definitions}

Let $\varphi\in \BV(\R^n)$, $s\in(0,1)$, and $p\in[1,\infty]$.  Following
\cite[Definition 2.1]{HTV}, a continuous path $X:[0,T]\to\R^n$ is
$(s,p)$-variable with respect to $\varphi$ if there is a relatively compact
open neighborhood $U$ of $X([0,T])$ such that
\[
 t\longmapsto \int_U |X_t-z|^{-n+1-s}\,\lVert D\varphi\rVert(dz)
 \quad\text{belongs to }L^p(0,T).
\]
For a matrix coefficient $\sigma=(\sigma_{jk})$, variability is required
componentwise; we write $X\in\V(\sigma,s,p)$.  The source abbreviation
$\V(\sigma,s)$ means $p=1$.

For a fixed starting point $x_0$ and driver $Y=(Y^1,Y^2)$, the Picard-type
integral map is
\[
 (\Tmap X)_t=x_0+\int_0^t \sigma(X_u)\,dY_u.
\]
In our example both paths are regular enough that this is simply the ordinary
Riemann--Stieltjes integral.

\section{Idea of the proof}

Variability measures how the occupation of a path interacts with the gradient
measure of the coefficient.  For a coefficient jumping across a line, a path
that leaves the line like $t\mapsto t^\alpha$ sees a potential of order
$t^{-\alpha s}$, which is integrable after choosing $\alpha sp<1$.  We choose
the two constant coefficient matrices on opposite sides of the line so that
they both send the driver velocity $(1,1)$ to the tangent vector $(1,0)$.
Consequently, although the input leaves the jump line transversally, the
integral process runs along it.  At a point of the line the relevant Riesz
potential integrates $|r|^{-1-s}$ against line measure and is infinite.

The matrix choice is arranged through the inverses.  Two positive-definite
inverse matrices share their first column $(1,1)$ and differ only in the
$(2,2)$ entry.  Sharing that column forces both original matrices to send
$(1,1)$ to $(1,0)$; varying only the inverse's $(2,2)$ entry also makes its rows
distributionally curl-free across a horizontal jump.

\section{The obstruction}

Set $L=\{(x_1,x_2)\in\R^2:x_2=0\}$ and define
\[
 B_+=\begin{pmatrix}1&1\\[2pt]1&2\end{pmatrix},
 \qquad
 B_-=\begin{pmatrix}1&1\\[2pt]1&3\end{pmatrix},
\]
\[
 A_+=B_+^{-1}=\begin{pmatrix}2&-1\\[2pt]-1&1\end{pmatrix},
 \qquad
 A_-=B_-^{-1}=\begin{pmatrix}\frac32&-\frac12\\[2pt]-\frac12&\frac12\end{pmatrix}.
\]
Let
\[
 \sigma(x_1,x_2)=
 \begin{cases}
 A_+,&x_2\geq0,\\
 A_-,&x_2<0.
 \end{cases}
\]
The value on $L$ is immaterial for the distributional derivative and for the
integral below.

\begin{theorem}[Jump-line failure of Picard invariance]\label{thm:main}
The coefficient $\sigma$ is bounded, belongs componentwise to
$\BV(\R^2)$, and is symmetric and uniformly positive definite.  It satisfies
the determinant, curl-free-inverse, and angle assumptions used in the source
paper's multidimensional Doss construction.

Fix $s\in(0,1)$ and $p\in[1,\infty)$.  Choose
\[
 0<\alpha<\min\left\{1,\frac1{sp}\right\},
 \qquad X_t=(0,t^\alpha),\qquad Y_t=(t,t),\quad 0\leq t\leq1.
\]
Then $X\in C^\alpha([0,1];\R^2)\cap\V(\sigma,s,p)$.  The driver is smooth
and both its coordinates are nonconstant, but with $x_0=0$,
\[
 (\Tmap X)_t=(t,0).
\]
Moreover,
\[
 \Tmap X\notin\V(\sigma,\rho,q)
 \quad\text{for every }\rho\in(0,1)\text{ and every }q\in[1,\infty].
\]
Thus the raw integral map does not preserve even the union of positive
variability classes.  In the basic $p=1$ setting one may take $X_t=(0,t)$, a
Lipschitz input, for every $s\in(0,1)$.
\end{theorem}

\begin{proof}
First, $\det A_+=1$ and $\det A_-=\frac12$.  The eigenvalues of $A_+$ are
$(3\pm\sqrt5)/2$, and those of $A_-$ are $1\pm1/\sqrt2$.  Hence
\[
 \left(1-\frac1{\sqrt2}\right)|\xi|^2
 \leq \langle \xi,\sigma(x)\xi\rangle
 \leq \frac{3+\sqrt5}{2}|\xi|^2
\]
for every $\xi$ and almost every $x$.  This proves bounded uniform positive
definiteness and the determinant condition.  It also gives the source angle
condition, since
\[
 \langle\xi,\sigma(x)\xi\rangle
 \geq \frac{1-1/\sqrt2}{(3+\sqrt5)/2}
       |\sigma(x)\xi|\,|\xi|.
\]

The coefficient is constant on the two half-planes.  Since
\[
 A_+-A_-=\frac12\begin{pmatrix}1&-1\\-1&1\end{pmatrix},
\]
each distributional derivative $D\sigma_{jk}$ is a nonzero signed constant
multiple of the normal vector measure on $L$.  In particular,
$\lVert D\sigma_{jk}\rVert=c_{jk}\HH^1\!\restriction_L$ with
$c_{jk}=1/2$.

The inverse field equals $B_+$ above $L$ and $B_-$ below it.  The first row is
the constant vector $(1,1)$.  The second row is $(1,b(x_2))$, where $b$ is a
one-dimensional step function.  Therefore both distributional row curls
vanish:
\[
 D_1(\sigma^{-1})_{k2}-D_2(\sigma^{-1})_{k1}=0,
 \qquad k=1,2.
\]
This is precisely the two-dimensional curl-free-inverse condition.

We next check input variability.  Choose any bounded open neighborhood $U$ of
$X([0,1])$.  Enlarging its line section only increases the potential, so for
$t>0$ and every entry,
\begin{align*}
 \int_U |X_t-z|^{-1-s}\,\lVert D\sigma_{jk}\rVert(dz)
 &\leq \frac12\int_{\R}(r^2+t^{2\alpha})^{-(1+s)/2}\,dr\\
 &=C_s t^{-\alpha s},
\end{align*}
where
$C_s=\frac12\int_{\R}(1+u^2)^{-(1+s)/2}\,du<\infty$.
The potential is infinite at $t=0$, but that single time is irrelevant to
$L^p(0,1)$.  Since $\alpha sp<1$,
\[
 \int_0^1 (C_s t^{-\alpha s})^p\,dt<\infty.
\]
Thus $X\in\V(\sigma,s,p)$ componentwise.

For almost every $t>0$, $X_t$ lies in the upper half-plane, so
$\sigma(X_t)=A_+$.  Since
\[
 A_+\binom11=\binom10,
\]
the ordinary Stieltjes integral gives
\[
 (\Tmap X)_t=\int_0^t A_+\binom11\,du=(t,0)=:Z_t.
\]

Finally, let $U$ be any relatively compact open neighborhood of the output
segment $Z([0,1])$.  For every $t\in[0,1]$, openness supplies
$\varepsilon_t>0$ such that
$\{(r,0):|r-t|<\varepsilon_t\}\subset U$.  Hence for every
$\rho\in(0,1)$ and every entry,
\[
 \int_U |Z_t-z|^{-1-\rho}\,\lVert D\sigma_{jk}\rVert(dz)
 \geq \frac12\int_{t-\varepsilon_t}^{t+\varepsilon_t}
 |t-r|^{-1-\rho}\,dr=+\infty.
\]
The potential is infinite at every time, so it belongs to no $L^q$ space.
This proves the claim.
\end{proof}

\begin{remark}[Compatibility with the source integral regime]
The Lipschitz path $Y$ belongs to $C^\gamma$ for every $\gamma<1$.  Given the
chosen $\alpha$ and $s$, take $\gamma\in(1-\alpha s,1)$.  Then
$\alpha s+\gamma>1$, exactly the source condition ensuring existence of the
Stieltjes integral for an $(s,1)$-variable input.  Here the integrand is in
fact constant almost everywhere, so the integral is classical.
\end{remark}

\section{A scalar positive subcase}

The obstruction is tangential: the output is forced to run inside a jump set.
The following elementary observation records a setting in which this cannot
happen.

\begin{proposition}[Positive scalar time driver]\label{prop:scalar}
Let $n=1$, $Y_t=t$, and $\sigma\in\BV(\R)\cap L^\infty(\R)$.  Suppose an input
path $X$ is such that the composition is well defined and, for some constants
$0<c\leq M<\infty$,
\[
 c\leq \sigma(X_t)\leq M\qquad\text{for almost every }t\in[0,T].
\]
Then
\[
 Z_t=x_0+\int_0^t\sigma(X_u)\,du
\]
belongs to $\V(\sigma,\rho)$ for every $\rho\in(0,1)$.
\end{proposition}

\begin{proof}
The path $Z$ is absolutely continuous and
$c(t-u)\leq Z_t-Z_u\leq M(t-u)$ for $u<t$.  Thus it is increasing and its
inverse on its image is $1/c$-Lipschitz.  Let $U$ be a relatively compact open
neighborhood of $Z([0,T])$ and put
$\mu=\lVert D\sigma\rVert\!\restriction_U$, a finite measure.  By Tonelli and
the monotone change of variables,
\begin{align*}
 \int_0^T\int_U |Z_t-z|^{-\rho}\,\mu(dz)\,dt
 &\leq \frac1c\int_U\int_{Z_0}^{Z_T}|y-z|^{-\rho}\,dy\,\mu(dz)\\
 &<\infty,
\end{align*}
because $\rho<1$, the two integration sets are bounded, and $\mu(U)<\infty$.
This is precisely $(\rho,1)$-variability.
\end{proof}

Proposition~\ref{prop:scalar} is not offered as a solution to the general open
problem.  It instead exposes the missing feature: a quantitative crossing
condition turns the occupation measure of the output into a controlled
one-dimensional measure.  In higher dimensions, uniform positive
definiteness alone does not prevent tangency to a codimension-one jump set, as
Theorem~\ref{thm:main} shows.

\section{Verification, limitations, and novelty}

The proof has no computational dependency.  Its checkable core consists of:
\begin{enumerate}
\item the displayed matrix inversions and eigenvalue bounds;
\item the identity $A_+(1,1)^\top=(1,0)^\top$;
\item the scaling
$\int_\R(r^2+a^2)^{-(1+s)/2}dr=C_sa^{-s}$;
\item divergence of $\int_{-\varepsilon}^{\varepsilon}|r|^{-1-\rho}dr$.
\end{enumerate}

The result does not assert that the source problem has a negative answer.  It
leaves open positive fixed-point results under output-transversality,
occupation-density, stochastic nonconcentration, or locally bounded gradient
potential assumptions.  The finite-$p$ strengthening also does not give an
$(s,\infty)$-variable input, because the common starting point lies on the
jump line.

A bounded search on 26 August 2026 covered the run's four lightweight indexes,
the exact source phrase and arXiv id, the authors' follow-ups
arXiv:2105.06249~\cite{HTV22} and arXiv:2407.06907~\cite{HTVrough}, and close
terms involving fixed points, integral-process variability, Riesz potentials,
and jump lines.  No matching statement was found.  Since the observation is
elementary and the search was not a systematic citation review, the correct
status is \emph{candidate substantial partial result, likely valid}, not a
definitive novelty claim.

\section{Human-review recommendation}

The packet is suitable for focused human review.  The most important semantic
question is whether the source authors intended the ``integral process'' as
the raw Picard image of arbitrary admissible inputs, as here, or only after
restricting the candidate set by an additional output-side occupation
condition.  Under the raw-map interpretation, the obstruction is exact.

\begin{thebibliography}{9}

\bibitem{HTV}
M. Hinz, J. M. T\"olle, and L. Viitasaari,
\emph{Variability of paths and differential equations with
$BV$-coefficients}, Ann. Inst. H. Poincar\'e Probab. Statist. 59 (2023),
2036--2082; \href{https://arxiv.org/abs/2003.11698}{arXiv:2003.11698}.

\bibitem{HTV22}
M. Hinz, J. M. T\"olle, and L. Viitasaari,
\emph{Sobolev regularity of occupation measures and paths, variability and
compositions}, Electron. J. Probab. 27 (2022), 1--29;
\href{https://arxiv.org/abs/2105.06249}{arXiv:2105.06249}.

\bibitem{HTVrough}
M. Hinz, J. M. T\"olle, and L. Viitasaari,
\emph{Variability and the existence of rough integrals with irregular
coefficients}, \href{https://arxiv.org/abs/2407.06907}{arXiv:2407.06907v2}
(2025).

\end{thebibliography}

\end{document}

